\chapter{1935:Frédéric Joliot and Irène Joliot-Curie}

\label{chap:chemistry1935}

The Nobel Prize in Chemistry for the year 1935 was awarded jointly to Frédéric Joliot and Irène Joliot-Curie.

\textbf{Motivation:}

in recognition of their synthesis of new radioactive elements

\section{Frédéric Joliot}

\subsection*{Early Life and Education}

Jean Frédéric Joliot, known as Frédéric Joliot, was born on 1900-03-19 in Paris, France.
He was a French physicist and chemist who shared the 1935 Nobel Prize in Chemistry with Irène Joliot-Curie for their synthesis of new radioactive elements.

\subsection*{Career and Major Research}

Joliot graduated in 1923 from the École de Physique et de Chimie Industrielles de la Ville de Paris.
In 1925 he became an assistant to Marie Curie at the Institut du Radium in Paris, and in 1926 he married Marie Curie's daughter Irène, with whom he thereafter worked closely.
In 1934 the couple discovered artificial radioactivity: they found that when aluminium is bombarded with alpha particles it is converted into a radioactive isotope of phosphorus, showing that radioactivity can be created in previously stable elements.
He became professor at the Collège de France in 1937.
After the Second World War he was appointed High Commissioner for Atomic Energy in 1945 and directed the creation of the French atomic energy programme until he was removed from the post in 1950.

\subsection*{Nobel Prize-winning Work}

The work for which Frédéric Joliot was awarded the Nobel Prize in Chemistry in 1935 was described by the Nobel Committee as: "in recognition of their synthesis of new radioactive elements".

Bombarding aluminium with alpha particles, the two laureates produced a radioactive form of phosphorus and showed that this new radioactivity did not pre-exist in the target material.
Stable elements could therefore be transformed into radioactive ones in the laboratory.

\subsection*{Significance and Impact}

The discovery of artificial radioactivity showed that radioactivity is not an exclusive property of naturally radioactive elements and provided scientists with an almost unlimited source of radioactive isotopes.
It became the basis of radioisotope chemistry, nuclear medicine, and much of later nuclear physics.

\subsection*{Later Career and Honors}

Frédéric Joliot was made Commander of the Legion of Honour and held honorary doctorates from several universities.
He died on 1958-08-14 in Paris, France.

\subsection*{Legacy}

The legacy of Frédéric Joliot endures through the lasting impact of his scientific contributions.
The artificial radioactivity that he and Irène Joliot-Curie discovered continues to be exploited throughout science and medicine.

\subsection*{Modern Developments}

The Joliot-Curies' discovery of artificial radioactivity opened the era of man-made radioisotopes, which are now produced in reactors and cyclotrons worldwide for nuclear medicine, medical imaging, cancer therapy, and research. Positron emission tomography (PET) relies directly on cyclotron-produced radiotracers descended from their work, and the Joliot-Curies' discovery also led to the discovery of nuclear fission.

\subsection*{Selected Publications}

\begin{itemize}

\item Joliot, F., \& Joliot-Curie, I. (1934). "Un nouveau type de radioactivité." Comptes rendus hebdomadaires des séances de l'Académie des sciences, 198, 254.

\item Curie, I., \& Joliot, F. (1932). "Émission de protons de grande énergie par les substances hydrogénées sous l'influence des rayons gamma très pénétrants." Comptes rendus hebdomadaires des séances de l'Académie des sciences, 194, 273.

\end{itemize}

\clearpage

\section{Irène Joliot-Curie}

\subsection*{Early Life and Education}

Irène Joliot-Curie, née Curie, was born on 1897-09-12 in Paris, France.
She was a French physicist and chemist who shared the 1935 Nobel Prize in Chemistry with Frédéric Joliot for their synthesis of new radioactive elements.

\subsection*{Career and Major Research}

The daughter of Pierre and Marie Curie, Irène was educated in Paris, and during the First World War she worked as a nurse-radiographer, helping to bring X-ray diagnosis to the battlefront.
From 1918 she worked as an assistant in the Curie laboratory at the Institut du Radium, and she received her doctorate in 1925 for a thesis on the alpha rays of polonium.
After her marriage to Frédéric Joliot in 1926, the two carried out their research together, and in 1934 they discovered artificial radioactivity.
She became professor at the Sorbonne in 1937, and in 1946 she succeeded her mother as director of the Institut du Radium.
In 1936 she briefly served as Under-Secretary of State for Scientific Research in the French government.

\subsection*{Nobel Prize-winning Work}

The work for which Irène Joliot-Curie was awarded the Nobel Prize in Chemistry in 1935 was described by the Nobel Committee as: "in recognition of their synthesis of new radioactive elements".

Together with her husband she showed that alpha particles can convert stable elements into radioactive isotopes, thereby synthesizing new radioactive elements in the laboratory.

\subsection*{Significance and Impact}

The discovery of artificial radioactivity, made by the two laureates, opened the way to the production of radioactive isotopes on demand.
It earned them the Nobel Prize in Chemistry in 1935, making them only the second married couple to share a Nobel Prize.

\subsection*{Later Career and Honors}

Irène Joliot-Curie directed the Institut du Radium from 1946 until her death.
She died on 1956-03-17 in Paris, France, of leukemia, an illness widely attributed to her lifelong exposure to radioactive substances.

\subsection*{Legacy}

The legacy of Irène Joliot-Curie endures through the lasting impact of her scientific contributions.
Her work, carried out together with her husband, remains fundamental to the use of radioactive isotopes in chemistry and medicine.

\subsection*{Modern Developments}

The discovery of artificial radioactivity by Irène and Frédéric Joliot-Curie opened the era of man-made radioisotopes, now produced in reactors and cyclotrons worldwide for nuclear medicine, medical imaging, cancer therapy, and research. Positron emission tomography (PET) relies directly on the cyclotron-produced radiotracers that descend from their work. Irène Joliot-Curie's contributions, together with her scientific legacy as the daughter of Marie Curie, helped establish the Curie family as a dynasty of modern nuclear science.

\subsection*{Selected Publications}

\begin{itemize}

\item Joliot, F., \& Joliot-Curie, I. (1934). "Un nouveau type de radioactivité." Comptes rendus hebdomadaires des séances de l'Académie des sciences, 198, 254.

\item Curie, I., \& Joliot, F. (1932). "Émission de protons de grande énergie par les substances hydrogénées sous l'influence des rayons gamma très pénétrants." Comptes rendus hebdomadaires des séances de l'Académie des sciences, 194, 273.

\end{itemize}

\clearpage

\section*{Chapter Summary}

The 1935 prize recognized Frédéric Joliot and Irène Joliot-Curie for their synthesis of new radioactive elements. Their discovery of artificial radioactivity extended radioactivity to elements produced in the laboratory and opened the way to radiotracer methods.
